\documentclass[10pt]{article}
\usepackage[T1]{fontenc}
\usepackage[utf8]{inputenc}
\usepackage[french]{babel}
\usepackage{amsmath,amssymb,mathrsfs,mathpazo}
\usepackage{geometry}
\geometry{a4paper,margin=2cm}
\usepackage{graphicx,array,tabularx,tabularray,colortbl,multirow,multicol,booktabs,enumitem,pifont,ulem,fancybox,diagbox,cancel}
\usepackage{pstricks,pstricks-add,pst-plot,pst-tree,pst-node,pst-text,pst-func,pst-math,pst-bspline,pst-3dplot,pst-eucl,pst-solides3d}
\let\pscalculate\undefined
\usepackage{pst-bezier}
\usepackage{tikz,pgfplots,tkz-tab}
\usetikzlibrary{arrows,arrows.meta,patterns,calc,positioning,trees,decorations.markings,shapes.geometric,shapes.misc,shapes.arrows,intersections,decorations.pathmorphing,decorations.pathreplacing,backgrounds,shadows,babel}
\pgfplotsset{compat=1.18}
\usepgfplotslibrary{fillbetween}
\usepackage[np]{numprint}
\usepackage{listings,makecell,pas-tableur,twoopt}
\let\barre\undefined
\usepackage[pstricks]{tablvar}
\definecolor{darkgreen}{RGB}{0,100,0}
\providecommand{\red}{\color{red}}
\providecommand{\blue}{\color{blue}}
\providecommand{\green}{\color{green}}
\providecommand{\black}{\color{black}}
\providecommand{\white}{\color{white}}
\providecommand{\gray}{\color{gray}}
\providecommand{\cyan}{\color{cyan}}
\providecommand{\magenta}{\color{magenta}}
\providecommand{\yellow}{\color{yellow}}
\providecommand{\violet}{\color{violet}}
\providecommand{\orange}{\color{orange}}
\providecommand{\euro}{\texteuro}
\providecommand{\EUR}{\texteuro}
\providecommand{\ieme}{\textsuperscript{e}}
\providecommand{\degres}{\ensuremath{^\circ}}
\providecommand{\degree}{\ensuremath{^\circ}}
\providecommand{\vv}[1]{\overrightarrow{#1}}
\providecommand{\Coord}[1]{#1}
\providecommand{\decosix}{}
\providecommand{\starredbullet}{\ensuremath{\bullet}}
\pagestyle{empty}
\setlength{\parindent}{0pt}
\usepackage{adjustbox}\DeclareUnicodeCharacter{2236}{:}\DeclareUnicodeCharacter{2190}{\ensuremath{\leftarrow}}
\providecommand{\R}{\mathbb{R}}
\providecommand{\N}{\mathbb{N}}
\providecommand{\Z}{\mathbb{Z}}
\providecommand{\Q}{\mathbb{Q}}
\providecommand{\D}{\mathbb{D}}
\providecommand{\C}{\mathbb{C}}
\providecommand{\e}{{\mathrm{e}}}
\providecommand{\vect}[1]{\overrightarrow{#1}}
\providecommand{\vectt}[1]{\overrightarrow{\mathrm{#1}}}
\providecommand{\vvt}[1]{\overrightarrow{\mathrm{#1}}}
\providecommand{\np}[2][]{{#2~#1}}
\providecommand{\vv}[1]{{\overrightarrow{#1}}}
\renewcommand{\d}{{\mathrm{d}}}
\providecommand{\degres}{{^\circ}}
\providecommand{\degre}{{^\circ}}
\providecommand{\blue}{\color{blue}}
\providecommand{\red}{\color{red}}
\providecommand{\enskip}{\quad}
\providecommand{\llbracket}{\lbrack\!\lbrack}
\providecommand{\rrbracket}{\rbrack\!\rbrack}
\providecommand{\euro}{{\text{€}}}
\providecommand{\no}{{\text{n°}}}
\providecommand{\interval}[2]{{[#1;#2]}}
\providecommand{\intervFO}[2]{{[#1;#2[}}
\providecommand{\dotfill}{\ldots}

\providecommand{\ds}{\displaystyle}
\providecommand{\textding}[1]{#1}
\providecommand{\Oij}{\(\mathrm{O};\vec{\imath},\vec{\jmath}\)}
\providecommand{\Oijk}{\(\mathrm{O};\vec{\imath},\vec{\jmath},\vec{k}\)}
\providecommand{\Ouv}{\(\mathrm{O};\vec{u},\vec{v}\)}
\providecommand{\intervFF}[2]{\ensuremath{[#1;#2]}}
\renewcommand{\euro}{\ifmmode\text{\texteuro}\else\texteuro\fi}
\DeclareUnicodeCharacter{00B0}{\ifmmode^{\circ}\else\textdegree\fi}
\DeclareUnicodeCharacter{20AC}{\ifmmode\text{\texteuro}\else\texteuro\fi}
% Énoncés : Éducation nationale. Transcription et corrigés : APMEP, Denis Vergès et contributeurs. Découpage et classement : CoopMaths, Sébastien Lozano, Éric Elter et contributeurs. Adaptation de la présentation : Maths Explorer.
% CC BY-SA : https://coopmaths.fr/www/about/
% https://coopmaths.fr/alea/static/e3c/2020/tex/e3c_2020_00_sujet17_2_cor.tex
\begin{document}
\noindent\textit{APMEP / CoopMaths. CC BY-SA. Présentation adaptée par Maths Explorer.}\par\medskip
On considère la fonction \(f\) définie sur \(\left] -1 ; +\infty \right[\) par :
	\[
	f(x) = \dfrac{x^2 + 1}{x + 1}.
	\]
	
	\subsection*{1.}
	
	Puisque \(x > -1\), alors \(x + 1 > 0\), \(f\) est donc une fonction quotient dérivable sur \(\left] -1 ; +\infty \right[\) et :
	\[
	f'(x) = \dfrac{2x(x + 1) - 1(x^2 + 1)}{(x + 1)^2} = \dfrac{2x^2 + 2x - x^2 - 1}{(x + 1)^2} = \dfrac{x^2 + 2x - 1}{(x + 1)^2}.
	\]
	
	\subsection*{2.}
	
	Comme, quel que soit le réel \(x \in \left] -1 ; +\infty \right[\), \((x + 1)^2 > 0\), le signe de \(f'(x)\) est celui du numérateur, le trinôme \(x^2 + 2x - 1\).
	
	Pour celui-ci \(\Delta = 4 + 4 = 4 \times 2 = (2\sqrt{2})^2\).\\ Le trinôme a donc deux racines :
	\[
	x_1 = \dfrac{-2 + 2\sqrt{2}}{2} = -1 + \sqrt{2}, \quad x_2 = \dfrac{-2 - 2\sqrt{2}}{2} = -1 - \sqrt{2}.
	\]
	
	On sait que le trinôme et donc la dérivée sont positifs, sauf sur l'intervalle \(\left] -1 - \sqrt{2} ; -1 + \sqrt{2} \right[\).
	
	Comme \(-1 - \sqrt{2} \approx -2,414\), la dérivée est négative sur \(\left] -1 ; -1 + \sqrt{2} \right[\) et positive sur \(\left] -1 + \sqrt{2} ; +\infty \right[\).
	
	La fonction est donc décroissante sur \(\left] -1 ; -1 + \sqrt{2} \right[\) et croissante sur \(\left] -1 + \sqrt{2} ; +\infty \right[\).
	
	\subsection*{3.}
	
	On sait que l'équation réduite de la tangente est : \(y - f(0) = f'(0)(x - 0)\).
	
	Avec \(f(0) = 1\) et \(f'(0) = -1\), on obtient :
	\[
	M(x ; y) \in Y \Leftrightarrow y - 1 = -x \Leftrightarrow y = -x + 1.
	\]
	
	\subsection*{4.}
	
	On étudie la fonction différence entre la fonction \(f\) et la fonction \(g\) telle que \(g(x) = x\), soit :
	\[
	d(x) = \dfrac{x^2 + 1}{x + 1} - x = \dfrac{x^2 + 1 - x(x + 1)}{x + 1} = \dfrac{x^2 + 1 - x^2 - x}{x + 1} = \dfrac{-x + 1}{x + 1}.
	\]
	
	Or le signe de ce quotient est le même que le signe du produit \((-x + 1)(x + 1)\).
	
	On sait que ce trinôme est négatif sauf entre ses racines \(-1\) et 1 où il est positif.
	
	Donc sur \(\left] -1 ; 1 \right[\), \(d(x) > 0\) : la courbe est au-dessus de la droite d'équation \(y = x\).
	
	Sur \(\left] 1 ; +\infty \right[\), \(d(x) < 0\) : la courbe est en dessous de la droite.

\end{document}
